%%%%%%%%%%  scan idx 108  |  folio 97  %%%%%%%%%%

\medskip
\noindent\textbf{Lema 4.2.}
Consideremos un diagrama
\[
\begin{tikzcd}[column sep=huge,row sep=large]
& R' \arrow[d] \\
R \arrow[r,"\mathrm{surjective}"'] & \bar R
\end{tikzcd}
\]
en $\underline C_\Lambda$, y sean $E$ y $E'$ módulos planos sobre
$R$ y $R'$, respectivamente. Entonces, para todo
$R'$-homomorfismo
\[
 E'\longrightarrow \bar E=E\otimes_R\bar R,
\]
el producto fibrado $E\times_{\bar E}E'$ es plano como módulo sobre
$R\times_{\bar R}R'$.

%%%%%%%%%%  scan idx 109  |  folio 98  %%%%%%%%%%

\medskip
\noindent\textit{Demostración.}
Sea exacta
\[
 0\longrightarrow I\longrightarrow R\longrightarrow\bar R
 \longrightarrow0
\]
de modo que
\[
 0\longrightarrow I\longrightarrow R\times_{\bar R}R'
 \longrightarrow R'\longrightarrow0
\]
es exacta. Como $E$ es plano sobre $R$,
\[
 0\longrightarrow I\otimes_RE\longrightarrow E
 \longrightarrow\bar E\longrightarrow0
\]
es exacta, y por tanto obtenemos la sucesión exacta
\begin{equation}\tag{*}
0\longrightarrow I\otimes_RE\longrightarrow
E\times_{\bar E}E'\longrightarrow E'\longrightarrow0.
\end{equation}
Pongamos $S=R\times_{\bar R}R'$. Como $I$ es nilpotente,
$E\times_{\bar E}E'$ es plano como $S$-módulo si y solo si
$(E\times_{\bar E}E')\otimes_SR'$ es plano como $R'$-módulo y
\[
 \operatorname{Tor}_1^S(E\times_{\bar E}E',R')=0.
\]
Pero
\[
 I\otimes_RE=I\otimes_S(E\times_{\bar E}E'),
 \qquad
 E'=R'\otimes_S(E\times_{\bar E}E'),
\]
y, por consiguiente, la sucesión exacta $(*)$ puede reescribirse como
\[
0\longrightarrow I\otimes_S(E\times_{\bar E}E')
\longrightarrow E\times_{\bar E}E'
\longrightarrow R'\otimes_S(E\times_{\bar E}E')
\longrightarrow0.
\]
Así,
\[
 \operatorname{Tor}_1^S(R',E\times_{\bar E}E')=0,
 \qquad
 R'\otimes_S(E\times_{\bar E}E')=E',
\]
y $E'$ es plano sobre $R'$. Se sigue que
$E\times_{\bar E}E'$ es plano sobre $S$. \hfill$\square$

\medskip
\noindent\textbf{Corolario 4.3.}
Dado un diagrama
\[
\begin{tikzcd}[column sep=huge,row sep=large]
(R,X) \arrow[dr,"{(\varphi,\Phi)}"']
&& (R',X') \arrow[dl,"{(\varphi',\Phi')}"] \\
& (\bar R,\bar X) &
\end{tikzcd}
\]
en $\underline M_{X_0}$, la suma amalgamada
$X\amalg_{\bar X}X'$ (en la categoría de esquemas) es un esquema
plano sobre $R\times_{\bar R}R'$, siempre que
$\varphi:R\longrightarrow\bar R$ sea sobreyectivo. En consecuencia,
$\underline M_{X_0}$ es una categoría cofibrada homogénea en
grupoides sobre $\underline C_\Lambda$.

\medskip
\noindent\textit{Demostración.}
Para todo abierto afín $U_0\subset X_0$, por definición tenemos
\[
 (X\amalg_{\bar X}X')|_{U_0}
 \simeq
 \operatorname{Spec}\!\left(
   \Gamma(U_0,X)
   \times_{\Gamma(U_0,\bar X)}
   \Gamma(U_0,X')
 \right).
\]
Por consiguiente, si $\varphi:R\longrightarrow\bar R$ es
sobreyectivo, $X\amalg_{\bar X}X'$ es plano sobre
$R\times_{\bar R}R'$ por el Lema 4.2. Además, el morfismo
\[
 i_X\amalg_{i_{\bar X}}i_{X'}:
 X_0\longrightarrow X\amalg_{\bar X}X'
\]
induce un isomorfismo
\[
 X_0\longrightarrow (X\amalg_{\bar X}X')\otimes_Sk,
 \qquad S=R\times_{\bar R}R'.
\]
Por tanto,
\[
 \bigl(R\times_{\bar R}R',X\amalg_{\bar X}X'\bigr)
 \in\operatorname{Ob}(\underline M_{X_0}),
\]
y, en $\underline M_{X_0}$, tenemos
\[
 \bigl(R\times_{\bar R}R',X\amalg_{\bar X}X'\bigr)
 =
 (R,X)\times_{(\bar R,\bar X)}(R',X').
\]
Así, $\underline M_{X_0}$ es homogénea sobre
$\underline C_\Lambda$. \hfill$\square$

%%%%%%%%%%  scan idx 110  |  folio 99  %%%%%%%%%%
